\documentclass[10pt,a4paper]{amsart}
\usepackage{amsmath,amssymb,amsthm,graphicx}
% Retained minimal preamble for the Stacks Project chapter flat.tex (More on Flatness).
% Verbatim excerpt of the chapter preamble: Xy-pic setup (native preamble lines 14-18) and the native
% theorem declarations with the mathematical macros (native preamble lines 154-217). Copied mechanically;
% nothing added, removed or redefined. The wrapper class and the amsmath/amssymb/amsthm/graphicx base
% packages are supplied by the assembler, so the class/externaldocument part of the native preamble is omitted.
% For commutative diagrams we use Xy-pic
\usepackage[all]{xy}

% We use 2cell for 2-commutative diagrams.
\xyoption{2cell}
% Theorem environments.
%
\theoremstyle{plain}
\newtheorem{theorem}[subsection]{Theorem}
\newtheorem{proposition}[subsection]{Proposition}
\newtheorem{lemma}[subsection]{Lemma}

\theoremstyle{definition}
\newtheorem{definition}[subsection]{Definition}
\newtheorem{example}[subsection]{Example}
\newtheorem{exercise}[subsection]{Exercise}
\newtheorem{situation}[subsection]{Situation}

\theoremstyle{remark}
\newtheorem{remark}[subsection]{Remark}
\newtheorem{remarks}[subsection]{Remarks}

\numberwithin{equation}{subsection}

% Macros
%
\def\lim{\mathop{\mathrm{lim}}\nolimits}
\def\colim{\mathop{\mathrm{colim}}\nolimits}
\def\Spec{\mathop{\mathrm{Spec}}}
\def\Hom{\mathop{\mathrm{Hom}}\nolimits}
\def\Ext{\mathop{\mathrm{Ext}}\nolimits}
\def\SheafHom{\mathop{\mathcal{H}\!\mathit{om}}\nolimits}
\def\SheafExt{\mathop{\mathcal{E}\!\mathit{xt}}\nolimits}
\def\Sch{\mathit{Sch}}
\def\Mor{\mathop{\mathrm{Mor}}\nolimits}
\def\Ob{\mathop{\mathrm{Ob}}\nolimits}
\def\Sh{\mathop{\mathit{Sh}}\nolimits}
\def\NL{\mathop{N\!L}\nolimits}
\def\CH{\mathop{\mathrm{CH}}\nolimits}
\def\proetale{{pro\text{-}\acute{e}tale}}
\def\etale{{\acute{e}tale}}
\def\QCoh{\mathit{QCoh}}
\def\Ker{\mathop{\mathrm{Ker}}}
\def\Im{\mathop{\mathrm{Im}}}
\def\Coker{\mathop{\mathrm{Coker}}}
\def\Coim{\mathop{\mathrm{Coim}}}

% Boxtimes
%
\DeclareMathSymbol{\boxtimes}{\mathbin}{AMSa}{"02}

%
% Macros for moduli stacks/spaces
%
\def\QCohstack{\mathcal{QC}\!\mathit{oh}}
\def\Cohstack{\mathcal{C}\!\mathit{oh}}
\def\Spacesstack{\mathcal{S}\!\mathit{paces}}
\def\Quotfunctor{\mathrm{Quot}}
\def\Hilbfunctor{\mathrm{Hilb}}
\def\Curvesstack{\mathcal{C}\!\mathit{urves}}
\def\Polarizedstack{\mathcal{P}\!\mathit{olarized}}
\def\Complexesstack{\mathcal{C}\!\mathit{omplexes}}
% \Pic is the operator that assigns to X its picard group, usage \Pic(X)
% \Picardstack_{X/B} denotes the Picard stack of X over B
% \Picardfunctor_{X/B} denotes the Picard functor of X over B
\def\Pic{\mathop{\mathrm{Pic}}\nolimits}
\def\Picardstack{\mathcal{P}\!\mathit{ic}}
\def\Picardfunctor{\mathrm{Pic}}
\def\Deformationcategory{\mathcal{D}\!\mathit{ef}}

% --- appended 2026-10-12 (Chief): the project own definitions of the reference/slogan/history
% display environments, copied mechanically from the shared Stacks preamble lines 8-12
% (\usepackage{verbatim} plus the three \newenvironment lines, which the chapter preamble keeps
% verbatim from the project preamble). Nothing added, removed or redefined: these environments
% carry no mathematical content and the project itself typesets them as comments. Without these
% three lines a counted statement whose own text uses them fails with Environment slogan undefined.
\usepackage{verbatim}
\newenvironment{reference}{\comment}{\endcomment}
\newenvironment{slogan}{\comment}{\endcomment}
\newenvironment{history}{\comment}{\endcomment}
\usepackage{cleveref}
\crefname{originallemma}{original lemma}{original lemmas}
\Crefname{originallemma}{Original lemma}{Original lemmas}
\makeatletter
\newlabel{lemma-affine-locally-projective-pure}{{05K6}{source}{}{sourcepointer-1}{}}
\newlabel{lemma-affine-locally-projective-pure@cref}{{[originallemma][0][]05K6}{[1][0][]source}{}{}{}}
\makeatother
\crefname{originallemma}{original lemma}{original lemmas}
\Crefname{originallemma}{Original lemma}{Original lemmas}
\makeatletter
\newlabel{lemma-associated-point-specializes}{{05L3}{source}{}{sourcepointer-2}{}}
\newlabel{lemma-associated-point-specializes@cref}{{[originallemma][0][]05L3}{[1][0][]source}{}{}{}}
\makeatother
\crefname{originallemma}{original lemma}{original lemmas}
\Crefname{originallemma}{Original lemma}{Original lemmas}
\makeatletter
\newlabel{lemma-etale-weak-assassin-up-down}{{05FP}{source}{}{sourcepointer-3}{}}
\newlabel{lemma-etale-weak-assassin-up-down@cref}{{[originallemma][0][]05FP}{[1][0][]source}{}{}{}}
\makeatother
\crefname{originallemma}{original lemma}{original lemmas}
\Crefname{originallemma}{Original lemma}{Original lemmas}
\makeatletter
\newlabel{lemma-finite-presentation-flat-along-fibre}{{05KW}{source}{}{sourcepointer-4}{}}
\newlabel{lemma-finite-presentation-flat-along-fibre@cref}{{[originallemma][0][]05KW}{[1][0][]source}{}{}{}}
\makeatother
\crefname{originallemma}{original lemma}{original lemmas}
\Crefname{originallemma}{Original lemma}{Original lemmas}
\makeatletter
\newlabel{lemma-flat-finite-presentation-purity-open}{{05MC}{source}{}{sourcepointer-5}{}}
\newlabel{lemma-flat-finite-presentation-purity-open@cref}{{[originallemma][0][]05MC}{[1][0][]source}{}{}{}}
\makeatother
\crefname{originallemma}{original lemma}{original lemmas}
\Crefname{originallemma}{Original lemma}{Original lemmas}
\makeatletter
\newlabel{more-morphisms-theorem-openness-flatness}{{0399}{source}{}{sourcepointer-6}{}}
\newlabel{more-morphisms-theorem-openness-flatness@cref}{{[originallemma][0][]0399}{[1][0][]source}{}{}{}}
\makeatother
\crefname{originallemma}{original lemma}{original lemmas}
\Crefname{originallemma}{Original lemma}{Original lemmas}
\makeatletter
\newlabel{proposition-finite-presentation-flat-pure-is-projective}{{05MD}{source}{}{sourcepointer-7}{}}
\newlabel{proposition-finite-presentation-flat-pure-is-projective@cref}{{[originallemma][0][]05MD}{[1][0][]source}{}{}{}}
\makeatother
\title{An affine etale neighbourhood on which the global sections of a flat finitely presented module are projective}
\author{The Stacks Project Authors}
\date{stacks-project:flat.tex:local-snapshot-20261011; GNU FDL 1.2 or later}
\begin{document}
\maketitle
\noindent\textit{Editorial note.} All delivered mathematical text is the author's own, copied from the pinned Stacks Project chapter file flat.tex; no statement or proof step was rewritten, reordered or removed.\par\medskip
\noindent\textit{Editorial note.} Preamble provenance: the chapter's own retained minimal preamble (sources/stacks-flat/retained-preamble.tex, lines 1-85) is a verbatim excerpt of the native Xy-pic setup and of the native theorem declarations and mathematical macros, followed by the project's own reference/slogan/history display-environment definitions copied mechanically from the shared preamble; the wrapper class and the amsmath/amssymb/amsthm/graphicx base packages are supplied by the assembler, and extra\_packages is left empty because everything needed is already in the retained preamble. No macro was invented or redefined.\par\medskip
\noindent\textit{Editorial note.} Same-chapter pointers are located at their \textbackslash{}label lines inside the pinned chapter file; cross-chapter pointers are located in the project's own tag list staged in the snapshot (sources/stacks-chow/tags.txt), at the verbatim lines carrying their canonical tags.\par\medskip
\noindent\textit{Allowed standard prerequisites (editorial).}\begin{itemize}
\item Openness of flatness. More on Morphisms, label more-morphisms-theorem-openness-flatness, invoked at proof line 5677. Lets the module be assumed flat over the whole of the shrunk source, which is the first reduction of the proof.
\item The free and projective variants of the flat along the fibre lemma, and purity of an affine locally projective module. Same chapter flat.tex, labels lemma-finite-presentation-flat-along-fibre (proof line 5684) and lemma-affine-locally-projective-pure (proof line 5696). Produces the affine etale neighbourhood and the etale g with X\_s inside g(X') and projective global sections, and gives universal purity of the pulled back module.
\item Purity of a flat finitely presented module on an open, the projective criterion, and specialization of associated points along an etale morphism of fibres. Same chapter flat.tex, labels lemma-flat-finite-presentation-purity-open (proof line 5703), proposition-finite-presentation-flat-pure-is-projective (proof line 5705), lemma-associated-point-specializes (proof line 5712) and lemma-etale-weak-assassin-up-down (proof line 5716). The first two upgrade purity of the restriction to projectivity of its global sections, and the others supply the specialization of an associated point of the restriction into the fibre, obtained by lifting the point along the etale morphism and using purity of the pulled back module.
\end{itemize}\medskip
\ifdefined\hypertarget\hypertarget{sourcepointer-1}{}\fi
\noindent\textit{Source reference.} Original lemma 05K6: sources/stacks-flat/pinned-original-source.tex, source lines 5148--5148. Navigation pointer for \textbackslash{}ref\{lemma-affine-locally-projective-pure\} used inside the counted span (line 5696); the argument that consumes it is retained in the span.\par\medskip
\ifdefined\hypertarget\hypertarget{sourcepointer-2}{}\fi
\noindent\textit{Source reference.} Original lemma 05L3: sources/stacks-flat/pinned-original-source.tex, source lines 5175--5175. Navigation pointer for \textbackslash{}ref\{lemma-associated-point-specializes\} used inside the counted span (line 5712); the argument that consumes it is retained in the span.\par\medskip
\ifdefined\hypertarget\hypertarget{sourcepointer-3}{}\fi
\noindent\textit{Source reference.} Original lemma 05FP: sources/stacks-flat/pinned-original-source.tex, source lines 353--353. Navigation pointer for \textbackslash{}ref\{lemma-etale-weak-assassin-up-down\} used inside the counted span (line 5716); the argument that consumes it is retained in the span.\par\medskip
\ifdefined\hypertarget\hypertarget{sourcepointer-4}{}\fi
\noindent\textit{Source reference.} Original lemma 05KW: sources/stacks-flat/pinned-original-source.tex, source lines 3342--3342. Navigation pointer for \textbackslash{}ref\{lemma-finite-presentation-flat-along-fibre\} used inside the counted span (line 5684); the argument that consumes it is retained in the span.\par\medskip
\ifdefined\hypertarget\hypertarget{sourcepointer-5}{}\fi
\noindent\textit{Source reference.} Original lemma 05MC: sources/stacks-flat/pinned-original-source.tex, source lines 5490--5490. Navigation pointer for \textbackslash{}ref\{lemma-flat-finite-presentation-purity-open\} used inside the counted span (line 5703); the argument that consumes it is retained in the span.\par\medskip
\ifdefined\hypertarget\hypertarget{sourcepointer-6}{}\fi
\noindent\textit{Source reference.} Original lemma 0399: sources/stacks-chow/tags.txt, source lines 4264--4264. Navigation pointer for \textbackslash{}ref\{more-morphisms-theorem-openness-flatness\} used inside the counted span (line 5677); the argument that consumes it is retained in the span.\par\medskip
\ifdefined\hypertarget\hypertarget{sourcepointer-7}{}\fi
\noindent\textit{Source reference.} Original lemma 05MD: sources/stacks-flat/pinned-original-source.tex, source lines 5582--5582. Navigation pointer for \textbackslash{}ref\{proposition-finite-presentation-flat-pure-is-projective\} used inside the counted span (line 5705); the argument that consumes it is retained in the span.\par\medskip
\begin{lemma}
\label{lemma-flat-finite-presentation-affine-neighbourhood-projective}
Let $f : X \to S$ be a morphism which is locally of finite presentation.
Let $\mathcal{F}$ be a quasi-coherent $\mathcal{O}_X$-module which is
of finite presentation. Let $x \in X$ with $s = f(x) \in S$.
If $\mathcal{F}$ is flat at $x$ over $S$ there exists an affine
elementary \'etale neighbourhood $(S', s') \to (S, s)$ and
an affine open $U' \subset X \times_S S'$ which contains $x' = (x, s')$
such that $\Gamma(U', \mathcal{F}|_{U'})$ is a projective
$\Gamma(S', \mathcal{O}_{S'})$-module.
\end{lemma}

\begin{proof}
During the proof we may replace $X$ by an open neighbourhood of $x$
and we may replace $S$ by an elementary \'etale neighbourhood of $s$.
Hence, by openness of flatness (see
More on Morphisms, Theorem \ref{more-morphisms-theorem-openness-flatness})
we may assume that $\mathcal{F}$ is flat over $S$.
We may assume $S$ and $X$ are affine.
After shrinking $X$ some more we may assume that any
point of $\text{Ass}_{X_s}(\mathcal{F}_s)$ is a generalization of $x$.
This property is preserved on replacing $(S, s)$ by an elementary
\'etale neighbourhood. Hence we may apply
Lemma \ref{lemma-finite-presentation-flat-along-fibre}
to arrive at the situation where there exists a diagram
$$
\xymatrix{
X \ar[dr] & & X' \ar[ll]^g \ar[ld] \\
& S &
}
$$
of schemes affine and of finite presentation over $S$,
where $g$ is \'etale, $X_s \subset g(X')$, and with
$\Gamma(X', g^*\mathcal{F})$ a projective $\Gamma(S, \mathcal{O}_S)$-module.
Note that in this case $g^*\mathcal{F}$ is universally pure over $S$, see
Lemma \ref{lemma-affine-locally-projective-pure}.

\medskip\noindent
Let $U \subset g(X')$ be an affine open neighbourhood of $x$.
We claim that $\mathcal{F}|_U$ is pure along $U_s$. If we prove this, then
the lemma follows because $\mathcal{F}|_U$ will be pure relative to $S$
after shrinking $S$, see
Lemma \ref{lemma-flat-finite-presentation-purity-open},
whereupon the projectivity follows from
Proposition \ref{proposition-finite-presentation-flat-pure-is-projective}.
To prove the claim we have to show, after replacing $(S, s)$
by an arbitrary elementary \'etale neighbourhood, that any point $\xi$ of
$\text{Ass}_{U/S}(\mathcal{F}|_U)$ lying over some
$s' \in S$, $s' \leadsto s$ specializes to a point of $U_s$.
Since $U \subset g(X')$ we can find a $\xi' \in X'$ with
$g(\xi') = \xi$. Because $g^*\mathcal{F}$ is pure over $S$, using
Lemma \ref{lemma-associated-point-specializes},
we see there exists a specialization $\xi' \leadsto x'$ with
$x' \in \text{Ass}_{X'_s}(g^*\mathcal{F}_s)$. Then
$g(x') \in \text{Ass}_{X_s}(\mathcal{F}_s)$ (see for example
Lemma \ref{lemma-etale-weak-assassin-up-down}
applied to the \'etale morphism $X'_s \to X_s$ of Noetherian schemes)
and hence $g(x') \leadsto x$ by our choice of $X$ above! Since
$x \in U$ we conclude that $g(x') \in U$. Thus
$\xi = g(\xi') \leadsto g(x') \in U_s$ as desired.
\end{proof}

\end{document}
